% External validation — DantinoX (JAX/Flax) vs. an independent PyTorch
% implementation (dllm's LLaDA), same architecture/loss/optimizer/data.
% Needs: graphicx, booktabs

\begin{figure}[t]
\centering
\includegraphics[width=0.48\textwidth]{cross_framework_validation.png}
\caption{%
  Training loss vs.\ tokens seen for DantinoX (JAX/Flax) and an independent
  PyTorch implementation (dllm's LLaDA), at matched architecture (512-d,
  8~heads, 12~layers, T5 SentencePiece vocabulary), matched masked-diffusion
  loss formula, and matched AdamW optimizer/schedule.
  Both converge from the theoretical initialization NLL
  ($\ln(\text{vocab})\approx10.4$) to ${\approx}4.6$ over the same
  262M~tokens.%
}
\label{fig:cross_framework_validation}
\end{figure}

\subsection{External Validation Against an Independent Implementation}
\label{sec:external_validation}

To validate that DantinoX's discrete-diffusion training is not an artefact
of a single codebase, we cross-check it against dllm~\cite{zhou2026dllm},
an independent, actively maintained PyTorch/HuggingFace-Trainer
implementation of LLaDA. We match, on both sides: architecture (512-d,
8~attention heads, 12~layers, gated SwiGLU-equivalent FFN with
expansion~4, RoPE, T5 SentencePiece vocabulary of 32{,}128, the same
reserved sentinel as the mask token); the exact masked-diffusion loss
formula ($t\sim\mathrm{Uniform}(0,1)$, linear masking schedule
$p_\mathrm{mask}=t$, per-token cross-entropy on masked positions divided
by $t\cdot L$); the optimizer (AdamW, identical learning rate, warmup,
and cosine schedule); and the training data (the same
WikiText-103-raw-v1 corpus, 50M-token cap, matched 90/10 train/validation
split). The two implementations differ only in framework
(JAX/Flax vs.\ PyTorch) and, necessarily, in random seeds and low-level
numerics.

Figure~\ref{fig:cross_framework_validation} shows both training- and
validation-loss trajectories over an identical token budget (262M
tokens, ${\approx}6$ epochs of the capped corpus). Both start from the
theoretical initialization loss ($\ln(32{,}128)\approx10.38$; observed
11.00 for DantinoX, 10.47 for dllm) and converge to a final train/val
loss of $4.591/4.670$ (DantinoX) and $4.554/4.548$ (dllm) --- within
$0.7\%$ of each other on the training loss after independently
re-implementing the objective, optimizer, and data pipeline in a
different language and framework. Neither shows a meaningful
train/validation gap at this scale, indicating neither is overfitting
the capped corpus.
We additionally compute a clean, unweighted evaluation perplexity
(masked-token cross-entropy averaged with fresh random maskings, i.e.\
without the $1/t$ training-time reweighting) on the held-out validation
split: DantinoX reaches $\mathrm{PPL}=182.2$ and dllm reaches
$\mathrm{PPL}=160.0$ --- the same order of magnitude, and consistent
with the small training-loss gap (perplexity amplifies small NLL
differences exponentially).

We repeated this comparison at a larger scale (768-d, 12~layers,
${\approx}138$M parameters) on the dllm side only, confirming a
consistent perplexity ($\mathrm{PPL}=133.8$, in line with the improvement
expected from more capacity under the same token budget). We were unable
to complete the matching DantinoX run at this scale: JAX/XLA's memory
usage grows across consecutive training steps for this configuration and
exhausts the 40GB accelerator regardless of batch size, optimizer choice,
or tensor-parallel sharding across 2 GPUs (which avoids the crash but
increases step time by over $30\times$ due to cross-device communication
at this model size). We were unable to identify the exact source of
this growth within our evaluation budget and report it as an open
engineering issue (Limitations) rather than omit the 768-d data point
entirely.

% ── Suggested addition to the Limitations section ──────────────────────────
% \paragraph{Cross-framework validation scope.} Our external validation
% against dllm (Section~\ref{sec:external_validation}) covers only the
% 512-d/12-layer, MHA, AdamW configuration end-to-end; the 768-d/12-layer
% point is validated on the dllm side only, since the matching DantinoX run
% hit a persistent, unresolved GPU memory growth issue across training
% steps at that scale. We did not find the underlying cause (ruled out:
% batch size, optimizer choice, tokenizer vocabulary size, and gradient
% checkpointing); 2-way tensor parallelism avoids the crash but is
% impractically slow at this model size (${>}30\times$ slower per step).
